\setcounterpageref{ex}{0}
\TaoPhieuKiemTra{11}

\setcounter{ex}{0}
\setcounter{bt}{0}

\textbf{\large I. PHẦN TRẮC NGHIỆM (3,0 điểm)}

{\itshape Khoanh tròn vào chữ cái đặt trước câu trả lời đúng. Mỗi câu đúng được 0,5 điểm.}

\begin{ex} % Câu 1
Số gồm $9$ chục và $0$ đơn vị viết là
\choice
{\True $90$}
{$9$}
{$900$}
{$19$}
\loigiai{
Số gồm $9$ chục và $0$ đơn vị là số tròn chục $90$.\\
\textbf{Đáp án đúng: A}
}
\end{ex}

\begin{ex} % Câu 2
Cho phép trừ có số bị trừ là $79$, số trừ là $15$. Hiệu là
\choice
{$54$}
{$74$}
{$94$}
{\True $64$}
\loigiai{
Hiệu là: $79 - 15 = 64$.\\
\textbf{Đáp án đúng: D}
}
\end{ex}

\begin{ex} % Câu 3
Một chiếc xô chứa được $10\text{ }l$ nước. Đã có sẵn $6\text{ }l$ trong xô. Cần đổ thêm bao nhiêu lít nước để đầy xô?
\choice
{$5\text{ }l$}
{\True $4\text{ }l$}
{$16\text{ }l$}
{$3\text{ }l$}
\loigiai{
Cần đổ thêm: $10 - 6 = 4\text{ }l$.\\
\textbf{Đáp án đúng: B}
}
\end{ex}

\begin{ex} % Câu 4
Cho điểm $B$ nằm giữa hai điểm $A$ và $C$ thẳng hàng. Điểm chung của hai đoạn thẳng $AB$ và $BC$ là
\choice
{Điểm $A$}
{Điểm $C$}
{\True Điểm $B$}
{Không có điểm chung}
\loigiai{
Hai đoạn thẳng $AB$ và $BC$ chung nhau điểm $B$.\\
\textbf{Đáp án đúng: C}
}
\end{ex}

\begin{ex} % Câu 5
So sánh: $18 - 9 \dots 15 - 7$.
\choice
{\True $>$}
{$<$}
{$=$}
{Không so sánh được}
\loigiai{
$18 - 9 = 9$; $15 - 7 = 8$. Vì $9 > 8$ nên điền dấu $>$.\\
\textbf{Đáp án đúng: A}
}
\end{ex}

\begin{ex} % Câu 6
Tính: $6 + 8 - 5 = \dots$
\choice
{$8$}
{$10$}
{$7$}
{\True $9$}
\loigiai{
$6 + 8 = 14$; $14 - 5 = 9$.\\
\textbf{Đáp án đúng: D}
}
\end{ex}

\vspace{4mm}
\textbf{\large II. PHẦN TỰ LUẬN (7,0 điểm)}

{\itshape Thí sinh trình bày chi tiết bài làm vào giấy kiểm tra.}

\begin{bt}[1,0 điểm]
Tính nhẩm:
\begin{multicols}{2}
\begin{enumerate}[a)]
  \item $9 + 5$
  \item $7 + 6$
  \item $13 - 8$
  \item $15 - 9$
\end{enumerate}
\end{multicols}
\loigiai{
\begin{itemize}
  \item $9 + 5 = 14$
  \item $7 + 6 = 13$
  \item $13 - 8 = 5$
  \item $15 - 9 = 6$
\end{itemize}
}
\end{bt}

\begin{bt}[1,5 điểm]
Đặt tính rồi tính:
\begin{multicols}{3}
\begin{enumerate}[a)]
  \item $45 + 40$
  \item $77 - 34$
  \item $63 + 54$
\end{enumerate}
\end{multicols}
\loigiai{
\begin{itemize}
  \item $45 + 40 = 85$
  \item $77 - 34 = 43$
  \item $63 + 54 = 117$
\end{itemize}
}
\end{bt}

\begin{bt}[1,0 điểm]
Tính giá trị biểu thức kèm đơn vị đo:
\begin{enumerate}[a)]
  \item $27\text{ kg} + 10\text{ kg} - 11\text{ kg}$
  \item $36\text{ }l - 12\text{ }l + 9\text{ }l$
\end{enumerate}
\loigiai{
\begin{enumerate}[a)]
  \item $27\text{ kg} + 10\text{ kg} - 11\text{ kg} = 37\text{ kg} - 11\text{ kg} = 26\text{ kg}$.
  \item $36\text{ }l - 12\text{ }l + 9\text{ }l = 24\text{ }l + 9\text{ }l = 33\text{ }l$.
\end{enumerate}
}
\end{bt}

\begin{bt}[1,5 điểm]
Bạn Lan có $26$ ngôi sao giấy. Bạn Mai có nhiều hơn bạn Lan $9$ ngôi sao giấy. Hỏi bạn Mai có bao nhiêu ngôi sao giấy?
\loigiai{
Số ngôi sao giấy bạn Mai có là:
\[ 26 + 9 = 35 \text{ (ngôi sao giấy)} \]
\textbf{Đáp số:} $35$ ngôi sao giấy.
}
\end{bt}

\begin{bt}[1,0 điểm]
Một thùng chứa đầy $46\text{ }\text{kg}$ đường kính. Buổi sáng cửa hàng đã bán ra $15\text{ }\text{kg}$. Hỏi trong thùng còn lại bao nhiêu ki-lô-gam đường kính?
\loigiai{
Trong thùng còn lại số ki-lô-gam đường kính là:
\[ 46 - 15 = 31\text{ }\text{kg} \]
\textbf{Đáp số:} $31\text{ }\text{kg}$ đường kính.
}
\end{bt}

\begin{bt}[1,0 điểm]
Tìm số bé nhất có hai chữ số mà tổng hai chữ số của nó bằng $9$.
\loigiai{
Để số bé nhất thì chữ số hàng chục bé nhất, chọn chữ số hàng chục là $1$.\\
Chữ số hàng đơn vị: $9 - 1 = 8$.\\
Vậy số đó là $18$.\\
\textbf{Đáp số:} $18$.
}
\end{bt}
